\section{中継機デバイス（Display）のプログラムソース}
本節では，中継機デバイスのモータ回転制御およびLEDマトリクス表示を行う制御プログラムの全文を示す．

\subsection{FerrisWheel.ino}
\begin{lstlisting}[caption={FerrisWheel.ino}, label={lst:disp_ino}]
#include "MotorControl.h"
#include "Display.h"
#include <WiFiS3.h>
#include <WiFiUdp.h>

const char SSID[] = "hackathon001-WPA2";
const char PASS[] = "hackathon001";

IPAddress localIP(192, 168, 11, 10);
IPAddress gateway(192, 168, 11, 1);
IPAddress subnet(255, 255, 255, 0);

const uint16_t LOCAL_PORT = 5005;

MotorControl motor;

WiFiUDP Udp;
uint8_t packetBuffer[10];

// 演奏中フラグ。'S'でtrue / 'E'でfalse。
// 'B'(BPM変更)は演奏中のみ受け付け、開始前の不意の'B'でモーターが回り出すのを防ぐ。
bool running = false;

void setup() {
    Serial.begin(115200);
    delay(2000);

    Serial.println("--- FerrisWheel Setup Start ---");

    motor.begin();
    motor.createRpmTable();

    WiFi.config(localIP, gateway, subnet);
    WiFi.begin(SSID, PASS);
    Serial.print("Connecting to WiFi...");
    while (WiFi.status() != WL_CONNECTED) {
        delay(500);
        Serial.print(".");
    }
    Serial.println("\nWiFi Connected!");
    Serial.print("IP: "); Serial.println(WiFi.localIP());

    Udp.begin(LOCAL_PORT);
    Serial.print("Listening on Unicast Port "); Serial.println(LOCAL_PORT);

    displaySetup();
}

void loop() {
    uint8_t packetSize = Udp.parsePacket();
    if (packetSize <= 0) return;

    if (packetSize != 2) {
        Serial.print("想定外のパケットサイズ: "); Serial.println(packetSize);
        while (Udp.available()) Udp.read();
        return;
    }

    Udp.read(packetBuffer, sizeof(packetBuffer));
    char    header  = packetBuffer[0];
    uint8_t payload = packetBuffer[1];

    switch (header) {
        case 'S':
            Serial.print("演奏開始 ('S') payload="); Serial.println(payload);
            // payload に BPM 段階が入っているのでそれで起動（元コードの global 変数参照バグを修正）
            running = true;
            motor.startRamp(payload);
            displayBPM(payload);
            // 冗長送信の残りパケットを捨てる
            while (Udp.parsePacket() > 0) {
                while (Udp.available()) Udp.read();
            }
            break;

        case 'B':
            if (!running) {
                Serial.println("演奏前の 'B' を無視");
                break;
            }
            Serial.print("BPM変更 ('B') 段階="); Serial.println(payload);
            motor.changeBPM(payload);
            displayBPM(payload);
            break;

        case 'E':
            Serial.println("演奏終了 ('E')");
            running = false;
            motor.stopRamp();
            clearDisplay();
            break;

        default:
            Serial.print("想定外のヘッダ: "); Serial.println(header);
            break;
    }
}
\end{lstlisting}

\subsection{Display.h}
\begin{lstlisting}[caption={Display.h}, label={lst:disp_h}]
#ifndef DISPLAY_H
#define DISPLAY_H

#include <Arduino.h>
#include "Arduino_LED_Matrix.h"

#define BPM_STEP_1  60
#define BPM_STEP_2  90
#define BPM_STEP_3  120
#define BPM_STEP_4  150
#define BPM_STEP_5  180

#define FONT_WIDTH   3
#define FONT_HEIGHT  5
#define MATRIX_ROWS  8
#define MATRIX_COLS  12

void displaySetup();
void displayBPM(uint8_t stepNumber);
void clearDisplay();
void drawDigit(int digit, int x_offset);

extern const uint8_t font3x5[10][15];
extern ArduinoLEDMatrix matrix;

#endif
\end{lstlisting}

\subsection{Display.cpp}
\begin{lstlisting}[caption={Display.cpp}, label={lst:disp_cpp}]
#include "Display.h"
#include "MotorControl.h"
#include "Arduino_LED_Matrix.h"

const uint8_t font3x5[10][15] = {
  // 0
  {1,1,1, 1,0,1, 1,0,1, 1,0,1, 1,1,1},
  // 1
  {0,1,0, 1,1,0, 0,1,0, 0,1,0, 1,1,1},
  // 2
  {1,1,1, 0,0,1, 1,1,1, 1,0,0, 1,1,1},
  // 3
  {1,1,1, 0,0,1, 0,1,1, 0,0,1, 1,1,1},
  // 4
  {1,0,1, 1,0,1, 1,1,1, 0,0,1, 0,0,1},
  // 5
  {1,1,1, 1,0,0, 1,1,1, 0,0,1, 1,1,1},
  // 6
  {1,1,1, 1,0,0, 1,1,1, 1,0,1, 1,1,1},
  // 7
  {1,1,1, 0,0,1, 0,0,1, 0,0,1, 0,0,1},
  // 8
  {1,1,1, 1,0,1, 1,1,1, 1,0,1, 1,1,1},
  // 9
  {1,1,1, 1,0,1, 1,1,1, 0,0,1, 1,1,1}
};

static const int bpmTable[5] = {
  BPM_STEP_1, BPM_STEP_2, BPM_STEP_3, BPM_STEP_4, BPM_STEP_5
};

ArduinoLEDMatrix matrix;

void displaySetup() {
    matrix.begin();
    uint32_t warmUp[3] = {0, 0, 0};
    matrix.loadFrame(warmUp);
    delay(200);
}

static uint8_t frameBuffer[MATRIX_ROWS][MATRIX_COLS];
static int currentBPM = -1;

void displayBPM(uint8_t stepNumber) {
    if (stepNumber < 1 || stepNumber > 5) return;
    uint8_t bpm = bpmTable[stepNumber - 1];

    memset(frameBuffer, 0, sizeof(frameBuffer));

    int digits[3];
    int numDigits;
    if (bpm >= 100) {
        numDigits = 3;
        digits[0] = bpm / 100;
        digits[1] = (bpm / 10) % 10;
        digits[2] = bpm % 10;
    } else if (bpm >= 10) {
        numDigits = 2;
        digits[0] = bpm / 10;
        digits[1] = bpm % 10;
    } else {
        numDigits = 1;
        digits[0] = bpm;
    }

    int totalWidth = numDigits * FONT_WIDTH + (numDigits - 1) * 1;
    int startX     = (MATRIX_COLS - totalWidth) / 2;

    for (int i = 0; i < numDigits; i++) {
        drawDigit(digits[i], startX + i * (FONT_WIDTH + 1));
    }

    uint32_t bitmap[3] = {0, 0, 0};
    for (int row = 0; row < MATRIX_ROWS; row++) {
        for (int col = 0; col < MATRIX_COLS; col++) {
            if (frameBuffer[row][col]) {
                int n = row * MATRIX_COLS + col;
                bitmap[n / 32] |= (1UL << (31 - (n % 32)));
            }
        }
    }
    matrix.loadFrame(bitmap);
}

void drawDigit(int digit, int x_offset) {
    if (digit < 0 || digit > 9) return;
    const int y_offset = (MATRIX_ROWS - FONT_HEIGHT) / 2;
    for (int y = 0; y < FONT_HEIGHT; y++) {
        for (int x = 0; x < FONT_WIDTH; x++) {
            int col = x_offset + x;
            int row = y_offset + y;
            if (col >= 0 && col < MATRIX_COLS && row >= 0 && row < MATRIX_ROWS) {
                frameBuffer[row][col] = font3x5[digit][y * FONT_WIDTH + x];
            }
        }
    }
}

void clearDisplay() {
    currentBPM = -1;
    uint32_t blank[3] = {0, 0, 0};
    matrix.loadFrame(blank);
}
\end{lstlisting}

\subsection{MotorControl.h}
\begin{lstlisting}[caption={MotorControl.h}, label={lst:motor_h}]
#ifndef MOTOR_CONTROL_H
#define MOTOR_CONTROL_H

#include <Arduino.h>

class MotorControl {
public:
    MotorControl();
    void begin();
    void createRpmTable();
    void startRamp(uint8_t stepNumber);
    void changeBPM(uint8_t stepNumber);
    void stopRamp();

private:
    void stepToNumber(uint8_t stepNumber);
    void updateStatus();

    const uint8_t PIN_MOTOR_PWM = 5;
    const uint8_t bpmTable[5]  = {60, 90, 120, 150, 180};
    const float   VCC           = 5.0;
    const uint8_t pwmTable[5]  = {26, 28, 30, 32, 37}; // 仮の数値

    uint8_t currentBPM = 0;
    float   RPM        = 0;
    float   omega      = 0;
    uint8_t pwmValue   = 0;
    float   V_average  = 0;
    float   rpmTable[5] = {};
};

#endif
\end{lstlisting}

\subsection{MotorControl.cpp}
\begin{lstlisting}[caption={MotorControl.cpp}, label={lst:motor_cpp}]
#include "MotorControl.h"

MotorControl::MotorControl() {}

void MotorControl::begin() {
    pinMode(PIN_MOTOR_PWM, OUTPUT);
    analogWrite(PIN_MOTOR_PWM, 0);
}

void MotorControl::createRpmTable() {
    for (uint8_t i = 0; i < 5; i++) {
        rpmTable[i] = bpmTable[i] / 16.0f;
    }
    Serial.println("[Motor] rpmTable 初期化完了");
}

void MotorControl::startRamp(uint8_t stepNumber) {
    stepToNumber(stepNumber);
    Serial.println("[Motor] Start Ramp Up...");
    for (int i = 0; i <= pwmValue; i += 5) {
        analogWrite(PIN_MOTOR_PWM, i);
        delay(40);
    }
    analogWrite(PIN_MOTOR_PWM, pwmValue);
    Serial.println("[Motor] Target speed reached.");
}

void MotorControl::stopRamp() {
    Serial.println("[Motor] Start Ramp Down...");
    for (int i = pwmValue; i >= 0; i -= 5) {
        analogWrite(PIN_MOTOR_PWM, i);
        delay(40);
    }
    analogWrite(PIN_MOTOR_PWM, 0);
    Serial.println("[Motor] Stopped.");
}

void MotorControl::changeBPM(uint8_t stepNumber) {
    stepToNumber(stepNumber);
    updateStatus();
    analogWrite(PIN_MOTOR_PWM, pwmValue);
    Serial.print("[Motor] Speed Changed. PWM: "); Serial.println(pwmValue);
}

void MotorControl::stepToNumber(uint8_t stepNumber) {
    if (stepNumber < 1 || stepNumber > 5) return;
    currentBPM = bpmTable[stepNumber - 1];
    pwmValue   = pwmTable[stepNumber - 1];
    RPM        = rpmTable[stepNumber - 1];
}

void MotorControl::updateStatus() {
    omega     = 2.0 * PI * RPM / 60.0f;
    V_average = VCC * (pwmValue / 256.0f);
}
\end{lstlisting}
